\documentclass[11pt]{article}

\usepackage[margin=1in]{geometry}
\usepackage{amsmath, amssymb, amsthm}
\usepackage{booktabs}
\usepackage{array}
\usepackage{listings}
\usepackage{xcolor}
\usepackage{hyperref}
\usepackage{enumitem}
\usepackage{microtype}

\hypersetup{
    colorlinks=true,
    linkcolor=blue!60!black,
    urlcolor=blue!60!black,
}

% SQL-ish listing style
\definecolor{kw}{RGB}{0,0,180}
\definecolor{cmt}{RGB}{100,100,100}
\definecolor{str}{RGB}{160,0,0}
\lstdefinelanguage{DECIQL}{
    keywords={SELECT,FROM,WHERE,DECIDE,IS,BOOLEAN,INTEGER,REAL,SUCH,THAT,
              AND,OR,NOT,WHEN,PER,SUM,MIN,MAX,COUNT,AVG,
              MINIMIZE,MAXIMIZE,CASE,THEN,ELSE,END},
    sensitive=true,
    morecomment=[l]{--},
    morestring=[b]',
}
\lstset{
    language=DECIQL,
    basicstyle=\ttfamily\small,
    keywordstyle=\color{kw}\bfseries,
    commentstyle=\color{cmt}\itshape,
    stringstyle=\color{str},
    showstringspaces=false,
    breaklines=true,
    columns=fullflexible,
    frame=single,
    framesep=4pt,
    xleftmargin=8pt,
    aboveskip=10pt,
    belowskip=10pt,
}

\title{Evaluation Order of \texttt{WHEN} and \texttt{PER} \\ in Aggregate Constraints}
\author{}
\date{}

\begin{document}
\maketitle

\section{Question}

When \texttt{WHEN} and \texttt{PER} compose on an aggregate constraint ---
\texttt{AGG(f(x)) op C WHEN W PER p} --- there are two possible evaluation
orders: filter rows first then group (\textbf{WHEN-PER}), or group first then
filter within each group (\textbf{PER-WHEN}). The following document formally
defines both orderings, characterizes exactly when they produce different
results, and analyzes the rewriting options available under each.

\section{Definitions}

\paragraph{Setup.} Table $T$ with rows indexed by $i$. Decision variables
$x_i$. Column $p : T \to G$ assigns each row to a group. Predicate
$W : T \to \{\text{true}, \text{false}\}$ is the \texttt{WHEN} condition. The
constraint has the form:
\[
\text{AGG}(f(x_i)) \diamond C \quad \text{WHEN } W \text{ PER } p
\]
where $\diamond \in \{=, \leq, \geq, <, >\}$.

\bigskip

\noindent\textbf{WHEN-PER} (filter, then group):
\[
T' \leftarrow \{i \in T : W(i)\}
\]
\[
\text{For each } g \in \{p(i) : i \in T'\}: \quad
\text{emit } \text{AGG}(\{f(x_i) : i \in T',\, p(i) = g\}) \diamond C
\]
Groups are drawn from the filtered set. A group exists only if it contains at
least one row passing $W$.

\medskip

\noindent\textbf{PER-WHEN} (group, then filter):
\[
\text{For each } g \in \{p(i) : i \in T\}: \quad
S_g \leftarrow \{i \in T : p(i) = g \wedge W(i)\}; \quad
\text{emit } \text{AGG}(S_g) \diamond C
\]
Groups are drawn from the full table. Every group emits a constraint; $W$ only
determines which rows contribute to the aggregate.

\bigskip

\paragraph{Observation.} For any group $g$, the contributing row set is
identical under both orderings: $\{i : p(i) = g \wedge W(i)\}$. The orderings
differ only in \emph{which groups emit constraints}:
\begin{itemize}
    \item WHEN-PER: $G_{\text{filtered}} = \{p(i) : W(i)\}$ --- groups with $\geq 1$ qualifying row
    \item PER-WHEN: $G_{\text{full}} = \{p(i) : i \in T\}$ --- all groups
\end{itemize}
Define $E = G_{\text{full}} \setminus G_{\text{filtered}}$: the
\textbf{empty groups}, where every row fails $W$.

\section{Divergence Condition}

For groups in $G_{\text{filtered}}$, both orderings emit the same constraint over the same rows.\\
For groups in $E$:
\begin{itemize}
    \item WHEN-PER: no constraint (group does not exist)
    \item PER-WHEN: $\text{AGG}(\emptyset) \diamond C$
\end{itemize}
\textbf{The orderings diverge if and only if $\text{AGG}(\emptyset) \diamond C$
is not vacuously true.}

\section{Case Analysis}

Under PER-WHEN, an empty group emits $\text{AGG}(\emptyset) \, \diamond \, C$.
Under WHEN-PER, it emits nothing. The orderings diverge when
$\text{AGG}(\emptyset) \, \diamond \, C$ can be false. We consider each
aggregate type.

\subsection*{SUM and COUNT}

SUM and COUNT over an empty set evaluate to $0$. Whether this causes a problem
depends on the constraint direction:

\begin{itemize}[leftmargin=2em]
    \item $0 \leq 500$: true. Upper bounds are satisfied.
    \item $0 \geq 3$: false. Lower bounds are violated.
    \item $0 = 10$: false. Equalities are violated.
\end{itemize}
\textbf{Diverges for lower bounds ($\geq$) and equalities ($=$). Immune for
upper bounds ($\leq$).}\\
Note: this is assuming $C>0$. If $C<0$, the opposite holds.

\subsection*{AVG}

AVG over an empty set is $0/0$ --- undefined. SQL returns NULL. Whether
NULL $\, \diamond \, C$ is treated as ``no constraint'' or as ``false'' is a
design choice. If it is treated as no constraint, the orderings converge. If
it is treated as false, they diverge. \textbf{We treat NULL as no constraint
(orderings converge). If the user needs empty groups to cause infeasibility,
the existence constraint from Section~5.3 handles this explicitly.}

\subsection*{MIN and MAX}

The divergence depends on what the constraint \emph{means}:

\begin{itemize}[leftmargin=2em]
    \item $\text{MAX}(\text{hours}) \leq 12$ means ``nobody exceeds 12 hours.'' If there is nobody, this is trivially true.
    \item $\text{MAX}(\text{hours}) \geq 8$ means ``somebody works at least 8 hours.'' If there is nobody, this is impossible.
\end{itemize}
More generally:
\begin{center}
\begin{tabular}{>{\raggedright\arraybackslash}p{0.23\textwidth} >{\raggedright\arraybackslash}p{0.29\textwidth} >{\raggedright\arraybackslash}p{0.19\textwidth} >{\raggedright\arraybackslash}p{0.12\textwidth}}
\toprule
\textbf{Constraint} & \textbf{Meaning} & \textbf{Over $\emptyset$} & \textbf{Diverges?} \\
\midrule
$\text{MAX}(e) \leq C$ & every value is $\leq C$ & trivially true & No \\
$\text{MIN}(e) \geq C$ & every value is $\geq C$ & trivially true & No \\
$\text{MAX}(e) \geq C$ & some value is $\geq C$ & impossible & \textbf{Yes} \\
$\text{MIN}(e) \leq C$ & some value is $\leq C$ & impossible & \textbf{Yes} \\
$\text{MIN/MAX}(e) = C$ & some value equals $C$ & impossible & \textbf{Yes} \\
\bottomrule
\end{tabular}
\end{center}
Constraints that require ``every row'' to satisfy a bound are trivially true
over an empty set. Constraints that require ``some row'' to satisfy a bound are
impossible over an empty set.

\subsection*{Summary}

The orderings diverge in two situations, both requiring an empty group:

\begin{enumerate}[leftmargin=2em]
    \item \textbf{SUM/COUNT} with a lower bound or equality ($\geq$, $=$)
    \item \textbf{MIN/MAX} with an existential meaning (``some row must satisfy \ldots'')
\end{enumerate}

Upper-bound SUM/COUNT and universal MIN/MAX (``every row must satisfy
\ldots'') are immune.

\section{Examples}


\subsection{Renewable Energy Standard (SUM $\geq$, diverges)}

Many jurisdictions mandate that each service region sources a minimum amount of
electricity from renewables. Grid operators dispatch power plants to meet
demand at minimum cost, subject to this renewable floor.

The constraint of interest: each region must generate at least 30\,MW from
renewable plants.

\begin{lstlisting}
SELECT p.plant_id, p.region, output_mw
FROM plants p
DECIDE output_mw IS REAL
SUCH THAT
    SUM(output_mw) >= 100 PER p.region
    AND SUM(output_mw) >= 30 WHEN p.is_renewable PER p.region
    AND output_mw <= p.max_capacity
MINIMIZE SUM(output_mw * p.cost)
\end{lstlisting}

Three constraints: (1)~each region produces $\geq 100$\,MW total,
(2)~at least 30\,MW of that must come from renewable plants,
(3)~no plant exceeds its capacity.

\begin{center}
\begin{tabular}{llcrr}
\toprule
\textbf{plant} & \textbf{region} & \textbf{is\_renewable} & \textbf{max\_capacity} & \textbf{cost} \\
\midrule
P1 & North & 1 & 50  & 8 \\
P2 & North & 1 & 40  & 9 \\
P3 & North & 0 & 200 & 3 \\
P4 & South & 0 & 150 & 4 \\
P5 & South & 0 & 100 & 5 \\
\bottomrule
\end{tabular}
\end{center}

Focus on the second constraint: \texttt{SUM(output\_mw) >= 30 WHEN p.is\_renewable PER p.region}.

\paragraph{Under WHEN-PER:}
\begin{enumerate}[leftmargin=2em]
    \item \textbf{Filter.} $T' = \{P1, P2\}$ (rows where \texttt{is\_renewable} is true).
    \item \textbf{Group.} Group $T'$ by region: $\{\text{North}\}$.
    \item \textbf{Emit.} One constraint: $\text{output}_{P1} + \text{output}_{P2} \geq 30$.
\end{enumerate}
South has no rows in $T'$, so no constraint is generated for South.
South meets its 100\,MW demand entirely from P4 and P5 (fossil plants).
The renewable floor is silently skipped.

\paragraph{Under PER-WHEN:}
\begin{enumerate}[leftmargin=2em]
    \item \textbf{Group.} Group $T$ by region: $\{\text{North}, \text{South}\}$.
    \item \textbf{Filter within group.}
        North: $S = \{P1, P2\}$.
        South: $S = \emptyset$ (no renewable plants).
    \item \textbf{Emit.}
        North: $\text{output}_{P1} + \text{output}_{P2} \geq 30$.
        South: $\text{SUM}(\emptyset) = 0 \geq 30$ --- \textbf{infeasible}.
\end{enumerate}
The query reports that South cannot meet the renewable mandate with its
current plant fleet. The operator learns this \emph{before dispatch} and
can procure renewable capacity.

\subsubsection*{Rewriting between orderings}

Under WHEN-PER, can we rewrite the constraint to get PER-WHEN behavior?
Yes --- move the filter from \texttt{WHEN} into the coefficient:

\begin{lstlisting}
-- Original (WHEN-PER, misses South):
SUM(output_mw) >= 30 WHEN p.is_renewable PER p.region

-- Rewritten (no WHEN, all groups constrained):
SUM(output_mw * p.is_renewable) >= 30 PER p.region
\end{lstlisting}

Since \texttt{is\_renewable} is $0/1$, renewable plants contribute
\texttt{output\_mw $\times$ 1} and fossil plants contribute
\texttt{output\_mw $\times$ 0}. With no \texttt{WHEN}, PER covers all
regions. South gets $\text{output}_{P4} \times 0 + \text{output}_{P5}
\times 0 = 0 \geq 30$ --- infeasible, as desired.

If the condition is not a boolean column (e.g., \texttt{WHEN efficiency > 0.5}),
it must be converted to a $0/1$ indicator explicitly:

\begin{lstlisting}
SUM(output_mw * CASE WHEN efficiency > 0.5 THEN 1 ELSE 0 END) >= 30 PER p.region
\end{lstlisting}

Any condition can be reduced to a binary mask this way. The nature of the
condition does not matter --- \textbf{what matters is whether the aggregate can absorb
the resulting zeros. This is a property of the aggregate function.}

\subsubsection*{Why the coefficient trick works for SUM}

The coefficient trick replaces non-qualifying rows with 0. Whether this
is safe depends on whether 0 behaves like ``not being there'' --- i.e.,
whether 0 is the \textbf{identity element} of the aggregate. The identity
element $e$ of an operation is the value such that $a \oplus e = a$ for
all $a$:

\begin{itemize}[leftmargin=2em]
    \item \textbf{SUM}: identity is $0$ ($a + 0 = a$). Phantom 0 changes
        nothing --- the trick works.
    \item \textbf{MIN}: identity is $+\infty$ ($\min(a, +\infty) = a$).
        Phantom 0 $\neq +\infty$, and since $0 \leq$ any non-negative value,
        MIN picks it up as a false minimum --- the trick fails.
    \item \textbf{MAX}: identity is $-\infty$ ($\max(a, -\infty) = a$).
        Phantom 0 $\neq -\infty$, but with non-negative variables $0 \leq$
        all real values, so MAX ignores it --- the trick happens to work.
    \item \textbf{AVG}: no identity exists. Any extra row shifts the
        denominator --- the trick fails.
\end{itemize}

The precise test is not ``does 0 equal the identity?'' but \textbf{does
including a phantom 0 change the aggregate's value?} For SUM and MAX
(with non-negative variables), it does not. For MIN and AVG, it does.
Section~5.3 develops a universal workaround that handles all aggregates.

\subsubsection*{The reverse: can PER-WHEN replicate WHEN-PER?}

Under PER-WHEN, every group emits a constraint --- including empty groups.
To replicate WHEN-PER, we would need empty groups to emit a vacuously true
constraint instead of $0 \geq 30$ (which is false).

One idea: make the RHS conditional. If the RHS were 0 for groups with no
renewable plants, the constraint would become $0 \geq 0$ (true) instead
of $0 \geq 30$ (false):

\[
\text{SUM}(\text{output\_mw}) \geq
\begin{cases}
30 & \text{if region has renewables} \\
0  & \text{otherwise}
\end{cases}
\]

But ``region has renewables'' is a group-level aggregate property ---
it depends on whether \emph{any} row in the group satisfies the condition.
A per-row coefficient trick cannot express this: each row only knows its
own \texttt{is\_renewable} value, not whether other rows in its group
qualify. There is no reverse coefficient trick.

\subsection{Hazardous Material Cap (SUM $\leq$, converges)}

Fire codes and OSHA regulations cap the total weight of hazardous material
stored in any single warehouse. Logistics operators maximize the total value
of goods stored, subject to this cap.

The constraint of interest: hazardous weight per warehouse must not exceed
500\,kg.

\begin{lstlisting}
SELECT i.item_id, i.warehouse_id, store
FROM items i
DECIDE store IS BOOLEAN
SUCH THAT
    SUM(i.weight_kg * store) <= 10000 PER i.warehouse_id
    AND SUM(i.weight_kg * store) <= 500 WHEN i.is_hazardous PER i.warehouse_id
MAXIMIZE SUM(store * i.value)
\end{lstlisting}

Two constraints: (1)~total weight per warehouse $\leq 10{,}000$\,kg,
(2)~hazardous weight per warehouse $\leq 500$\,kg.

\begin{center}
\begin{tabular}{llcrr}
\toprule
\textbf{item} & \textbf{warehouse} & \textbf{is\_hazardous} & \textbf{weight\_kg} & \textbf{value} \\
\midrule
I1 & A & 1 & 120 & 800  \\
I2 & A & 1 & 200 & 600  \\
I3 & A & 0 & 500 & 400  \\
I4 & B & 0 & 300 & 900  \\
I5 & B & 0 & 450 & 700  \\
\bottomrule
\end{tabular}
\end{center}

Focus on the second constraint:
\texttt{SUM(i.weight\_kg * store) <= 500 WHEN i.is\_hazardous PER i.warehouse\_id}.

\paragraph{Under WHEN-PER:}
\begin{enumerate}[leftmargin=2em]
    \item \textbf{Filter.} $T' = \{I1, I2\}$ (rows where \texttt{is\_hazardous} is true).
    \item \textbf{Group.} Group $T'$ by warehouse: $\{A\}$.
    \item \textbf{Emit.} One constraint: $120 \cdot \text{store}_{I1} + 200 \cdot \text{store}_{I2} \leq 500$.
\end{enumerate}
Warehouse B has no rows in $T'$, so no constraint is generated for B.

\paragraph{Under PER-WHEN:}
\begin{enumerate}[leftmargin=2em]
    \item \textbf{Group.} Group $T$ by warehouse: $\{A, B\}$.
    \item \textbf{Filter within group.}
        A: $S = \{I1, I2\}$.
        B: $S = \emptyset$ (no hazardous items).
    \item \textbf{Emit.}
        A: $120 \cdot \text{store}_{I1} + 200 \cdot \text{store}_{I2} \leq 500$.
        B: $\text{SUM}(\emptyset) = 0 \leq 500$ --- \textbf{trivially true}.
\end{enumerate}

\textbf{Both orderings agree.} WHEN-PER generates no constraint for B;
PER-WHEN generates a vacuously true constraint for B. Neither restricts
warehouse B's variables. The feasible region is identical.

\subsubsection*{Why upper bounds converge}

In the energy example (Section~5.1), the lower bound $0 \geq 30$ was false,
causing divergence. Here the upper bound $0 \leq 500$ is true. The
difference comes from where 0 sits relative to the bound:

\begin{itemize}[leftmargin=2em]
    \item \textbf{Lower bounds} ($\geq C$, $C > 0$): the empty aggregate 0 is
        \emph{below} the threshold --- it fails.
    \item \textbf{Upper bounds} ($\leq C$, $C \geq 0$): the empty aggregate 0 is
        \emph{below} the cap --- it passes.
\end{itemize}

This is the same identity-element observation from Section~5.1: $\text{SUM}(\emptyset) = 0$.
For lower bounds, 0 is on the wrong side. For upper bounds, 0 is on the safe side.

\subsubsection*{Rewriting is unnecessary}

The coefficient trick still works --- \texttt{SUM(i.weight\_kg * store *
i.is\_hazardous) <= 500 PER i.warehouse\_id} --- but there is no reason to
use it. Both orderings already produce the same feasible region. Rewriting
only matters when the orderings diverge.

\subsubsection*{Edge case: negative upper bounds}

Convergence depends on $C \geq 0$. If the RHS were negative --- e.g.,
$\leq -5$ --- then $0 \leq -5$ is false and the orderings would diverge.
Negative upper bounds are rare in practice (they assert that an aggregate
must be negative), but they break the convergence guarantee.

\subsection{Certified Operator Requirement (MAX $\geq$, diverges)}

Safety regulations require each production line to have at least one operator
with an electrical or safety certification working $\geq 8$ hours per shift.
Operators may hold other certifications (e.g., mechanical), but only these
two qualify for the safety mandate. Plant managers minimize total labor cost
subject to this requirement.

The constraint of interest: on each line, the maximum hours worked by a
qualifying operator must be $\geq 8$.

\begin{lstlisting}
SELECT o.operator_id, o.line_id, o.certification, o.hourly_rate, hours
FROM operators o
DECIDE hours IS INTEGER
SUCH THAT
    SUM(hours) >= 16 PER o.line_id
    AND MAX(hours) >= 8
        WHEN (o.certification = 'electrical' OR o.certification = 'safety')
        PER o.line_id
    AND hours <= 12
MINIMIZE SUM(hours * o.hourly_rate)
\end{lstlisting}

Three constraints: (1)~each line has $\geq 16$ total operator-hours,
(2)~at least one electrical/safety-certified operator works $\geq 8$ hours,
(3)~no operator works more than 12 hours.

\begin{center}
\begin{tabular}{llcr}
\toprule
\textbf{operator} & \textbf{line} & \textbf{certification} & \textbf{hourly\_rate} \\
\midrule
Op1 & L1 & electrical  & \$35 \\
Op2 & L1 & mechanical  & \$28 \\
Op3 & L2 & mechanical  & \$28 \\
Op4 & L2 & mechanical  & \$25 \\
\bottomrule
\end{tabular}
\end{center}

Line L2 has certified operators --- but not with the right certification.
Both Op3 and Op4 hold mechanical certifications, which do not satisfy the
electrical/safety requirement.

Focus on the second constraint:
\texttt{MAX(hours) >= 8 WHEN (certification = 'electrical' OR certification = 'safety') PER line\_id}.

\paragraph{Under WHEN-PER:}
\begin{enumerate}[leftmargin=2em]
    \item \textbf{Filter.} $T' = \{\text{Op1}\}$ (only Op1 has a qualifying certification).
    \item \textbf{Group.} Group $T'$ by line: $\{L1\}$.
    \item \textbf{Emit.} One constraint: $\text{hours}_{\text{Op1}} \geq 8$.
\end{enumerate}
L2 has no rows in $T'$, so no constraint is generated. L2 operates with
only mechanically-certified staff --- the safety gap goes undetected.

\paragraph{Under PER-WHEN:}
\begin{enumerate}[leftmargin=2em]
    \item \textbf{Group.} Group $T$ by line: $\{L1, L2\}$.
    \item \textbf{Filter within group.}
        L1: $S = \{\text{Op1}\}$.
        L2: $S = \emptyset$ (no electrical/safety operators).
    \item \textbf{Emit.}
        L1: $\text{hours}_{\text{Op1}} \geq 8$.
        L2: $\text{MAX}(\emptyset) \geq 8$ --- no candidates exist --- \textbf{infeasible}.
\end{enumerate}
The query reports that L2 cannot satisfy the safety mandate. The plant
manager learns this before scheduling and can reassign a qualified operator.

\subsubsection*{Why this differs from SUM}

In the energy example (Section~5.1), the empty group produced $0 \geq 30$
--- a concrete falsehood. MAX over an empty group also produces a concrete
falsehood, just through different mechanics. \texttt{MAX(hours) >= 8} is a ``hard'' case --- it requires that
\emph{some} row satisfies the bound (existential), which cannot be
expressed as simple per-row constraints. It gets linearized with Big-M indicators:
one binary $y_i$ per row, with $\text{SUM}(y_i) \geq 1$ (``at least one
row must be selected''). For L2, there are no rows, so no $y_i$ variables
are created, and the constraint becomes $0 \geq 1$ --- infeasible.

The key difference from SUM is \emph{semantic}. MAX constraints have two
directions:

\begin{itemize}[leftmargin=2em]
    \item $\text{MAX}(\text{hours}) \geq 8$: ``\emph{some} operator works
        $\geq 8$ hours'' (existential). Over an empty set, this is
        impossible --- there is nobody to satisfy it.
    \item $\text{MAX}(\text{hours}) \leq 12$: ``\emph{every} operator works
        $\leq 12$ hours'' (universal). Over an empty set, this is vacuously
        true --- there is nobody to violate it.
\end{itemize}

If the constraint were \texttt{MAX(hours) <= 12 WHEN (\ldots) PER
line\_id}, L2 would pass. Same data, same empty group, opposite outcome.
The divergence depends on whether the constraint is existential (``some
row'') or universal (``every row'').

\subsubsection*{Rewriting: can WHEN-PER catch this?}

The universal workaround is an \textbf{existence constraint}: a guard that
guarantees every group has at least one qualifying row.

\begin{lstlisting}
-- Guard: every line must have at least one qualifying operator
SUM(CASE WHEN certification IN ('electrical', 'safety')
         THEN 1 ELSE 0 END) >= 1 PER o.line_id
-- Now safe: WHEN-PER will never see an empty group
AND MAX(hours) >= 8
    WHEN (certification = 'electrical' OR certification = 'safety')
    PER o.line_id
\end{lstlisting}

The guard has no \texttt{WHEN}, so PER covers all lines. A line with no
qualifying operators fails the guard ($0 \geq 1$), making the model
infeasible. With the guard in place, $E = \emptyset$: every group has
qualifying rows, so WHEN-PER generates constraints for all groups ---
identical to PER-WHEN.

\subsubsection*{Universality of the existence constraint}

The existence constraint works for \emph{every} diverging constraint case
in this document. By eliminating empty groups at the source, it makes the
ordering question irrelevant --- both orderings produce the same constraints
over the same row sets. The pattern is always the same:

\begin{center}
\texttt{SUM(CASE WHEN condition THEN 1 ELSE 0 END) >= 1 PER group\_col}
\end{center}

This works for SUM/COUNT with lower bounds (Section~5.1), existential MAX,
existential MIN --- any aggregate, any operator direction.

The \textbf{coefficient trick} from Section~5.1 is a useful optimization
that avoids the extra constraint: for SUM and MAX, embedding the condition
as a $0/1$ multiplier produces correct results because phantom 0 does not
distort these aggregates. For MIN and AVG, phantom 0 \emph{does} distort
the result, so the coefficient trick fails and the existence constraint is
the only option.

\subsubsection*{The reverse: can PER-WHEN replicate WHEN-PER?}

Under PER-WHEN, L2 emits $\text{MAX}(\emptyset) \geq 8$ --- infeasible.
To replicate WHEN-PER (where L2 gets no constraint), we would need to make
the constraint vacuously true for empty groups. The same barrier from
Section~5.1 applies: ``group has qualifying rows'' is a group-level
aggregate property, not a per-row column --- it cannot be expressed as a
conditional RHS within a single constraint. There is no reverse trick.

\subsection{Regional Revenue Equity (Objective: MAXIMIZE MIN(SUM) PER, diverges)}

A retailer expanding into new markets wants to maximize the worst region's
organic product revenue, ensuring balanced growth. This is a max-min fairness
objective: raise the floor across regions rather than concentrating on the
strongest one.

\begin{lstlisting}
SELECT p.product_id, p.region, p.is_organic, p.revenue, p.shelf_cost, stock
FROM products p
DECIDE stock IS BOOLEAN
SUCH THAT
    SUM(p.shelf_cost * stock) <= 10000 PER p.region
MAXIMIZE MIN(SUM(p.revenue * stock) PER p.region) WHEN p.is_organic
\end{lstlisting}

The objective has two levels. First, rows are grouped by region (PER),
and within each group the \textbf{inner} aggregate
$\text{SUM}(\text{revenue} \times \text{stock})$ computes that region's
organic revenue. Then the \textbf{outer} aggregate MIN selects the worst
region's total. WHEN controls which rows contribute to each group's inner
SUM. The ordering question --- WHEN-PER vs PER-WHEN --- determines whether
a region with no qualifying rows is excluded from the outer MIN entirely,
or included with a phantom 0.

\begin{center}
\begin{tabular}{llcrr}
\toprule
\textbf{product} & \textbf{region} & \textbf{is\_organic} & \textbf{revenue} & \textbf{shelf\_cost} \\
\midrule
P1  & West    & 1 & 500 & 2000 \\
P2  & West    & 1 & 300 & 1500 \\
P3  & East    & 1 & 400 & 1800 \\
P4  & East    & 1 & 600 & 2500 \\
P5  & Central & 0 & 800 & 3000 \\
P6  & Central & 0 & 700 & 2800 \\
\bottomrule
\end{tabular}
\end{center}

Central has products, but none are organic.

\paragraph{Under WHEN-PER:}
\begin{enumerate}[leftmargin=2em]
    \item \textbf{Filter.} $T' = \{P1, P2, P3, P4\}$ (organic products only).
    \item \textbf{Group.} Group $T'$ by region: $\{\text{West}, \text{East}\}$.
    \item \textbf{Inner SUM per group.}
        West: $500 \cdot \text{stock}_{P1} + 300 \cdot \text{stock}_{P2}$.
        East: $400 \cdot \text{stock}_{P3} + 600 \cdot \text{stock}_{P4}$.
    \item \textbf{Outer MIN.} $\text{MIN}(\text{West}_{\text{sum}},\, \text{East}_{\text{sum}})$.
\end{enumerate}
Central is excluded. The optimizer balances organic revenue between West and
East --- e.g., stocking products in both regions to raise the minimum.

\paragraph{Under PER-WHEN:}
\begin{enumerate}[leftmargin=2em]
    \item \textbf{Group.} Group $T$ by region: $\{\text{West}, \text{East}, \text{Central}\}$.
    \item \textbf{Filter within group.}
        West: $S = \{P1, P2\}$.
        East: $S = \{P3, P4\}$.
        Central: $S = \emptyset$.
    \item \textbf{Inner SUM per group.}
        West and East: same as above.
        Central: $\text{SUM}(\emptyset) = 0$.
    \item \textbf{Outer MIN.} $\text{MIN}(\text{West}_{\text{sum}},\, \text{East}_{\text{sum}},\, 0)$.
\end{enumerate}
Central contributes a phantom 0 to the outer MIN. Since
$0 \leq \text{West}_{\text{sum}}$ and $0 \leq \text{East}_{\text{sum}}$ for
any non-negative stocking, the MIN is pinned at 0. The optimizer cannot
increase it --- Central has no organic products to stock. \textbf{The
objective value is stuck at 0 regardless of what West and East achieve.}

\subsubsection*{What the inner aggregate produces for empty groups}

Under PER-WHEN, an empty group still participates in the inner aggregate.
What value does it contribute to the outer?

\begin{center}
\begin{tabular}{>{\raggedright\arraybackslash}p{0.18\textwidth}
                >{\raggedright\arraybackslash}p{0.50\textwidth}
                >{\raggedright\arraybackslash}p{0.18\textwidth}}
\toprule
\textbf{Inner agg.} & \textbf{Why} & \textbf{Enters outer as} \\
\midrule
SUM, COUNT, AVG &
    Sum of no rows is 0. COUNT reduces to SUM; AVG scales
    coefficients by $1/n_g$, but with no rows there is nothing
    to scale. &
    Phantom 0 \\[6pt]
MAX (easy) &
    The auxiliary $z_g$ must be $\geq$ every row's value (universal
    direction --- ``every row satisfies''). With no rows, nothing
    forces $z_g$ above 0. The optimizer pushes it to its lower
    bound. &
    identity ($-\infty$) \\[6pt]
MIN (easy) &
    The auxiliary $z_g$ must be $\leq$ every row's value (universal
    direction). With no rows, nothing caps $z_g$ from above. If the
    outer is MIN, $z_g$ never binds (harmless). If the outer is SUM,
    $z_g \to \infty$ (unbounded). &
    identity ($\infty$) \\[6pt]
MIN, MAX (hard) &
    Hard cases are the existential direction (``at least one row must
    achieve the bound''), requiring indicator variables with
    $\text{SUM}(y_i) \geq 1$. With no rows, there are no indicators:
    $0 \geq 1$. &
    \textbf{Infeasible} \\
\bottomrule
\end{tabular}
\end{center}

Inner SUM is the common case. The remaining cases (inner MIN/MAX) follow
the same easy/hard classification from Section~4. Min/Max inners seem very weird to have in the objective so for now just ignore for simplicity.

\subsubsection*{Does the phantom 0 distort the outer?}

The same identity element analysis from Section~5.1 applies. The outer
aggregate receives one value per group --- including the phantom 0 from
empty groups. Does that 0 change the outer's value?

\begin{center}
\begin{tabular}{>{\raggedright\arraybackslash}p{0.12\textwidth}
                >{\raggedright\arraybackslash}p{0.12\textwidth}
                >{\raggedright\arraybackslash}p{0.58\textwidth}}
\toprule
\textbf{Outer} & \textbf{Identity} & \textbf{Phantom 0 distorts?} \\
\midrule
SUM & $0$ & No --- $0$ is the identity \\
MAX & $-\infty$ & No --- $0 \leq$ real values, ignored \\
MIN & $+\infty$ & \textbf{Yes} --- $0 \leq$ real values, picked up as false minimum \\
AVG & (none) & No --- reduces to SUM$/N$; phantom 0 adds 0 to numerator
    (absorbed) and $N$ is a constant divisor (same argmax) \\
\bottomrule
\end{tabular}
\end{center}

Only outer MIN diverges. SUM and MAX absorb phantom 0 directly. AVG
reduces to SUM divided by a constant --- the denominator differs between
orderings but does not affect which solution is optimal.

In this example, the inner is SUM (produces phantom 0) and the outer is
MIN (distorted by phantom 0). Central's $\text{SUM}(\emptyset) = 0$ enters
the outer MIN, pinning the objective at 0.

\subsubsection*{Rewriting: can WHEN-PER replicate PER-WHEN?}

The coefficient trick embeds the WHEN condition as a $0/1$ multiplier in
the inner aggregate, removing the \texttt{WHEN} clause so all groups
participate. Whether this correctly replicates PER-WHEN depends on whether
the inner aggregate absorbs phantom 0 --- the same identity element test
from Section~5.1, applied at the inner level:

\begin{itemize}[leftmargin=2em]
    \item \textbf{Inner SUM, COUNT}: phantom 0 absorbed. Works.
    \item \textbf{Inner MAX}: phantom 0 ignored ($0 \leq$ real values). Works.
    \item \textbf{Inner MIN}: phantom 0 picked up as false minimum. Fails.
    \item \textbf{Inner AVG}: extra zeroed rows shift the denominator. Fails.
\end{itemize}

For this example (inner SUM):

\begin{lstlisting}
-- Original (WHEN-PER, excludes Central):
MAXIMIZE MIN(SUM(p.revenue * stock) PER p.region) WHEN p.is_organic

-- Rewritten (no WHEN, all regions participate):
MAXIMIZE MIN(SUM(p.revenue * stock * p.is_organic) PER p.region)
\end{lstlisting}

Central's rows contribute $\text{revenue} \times \text{stock} \times 0 = 0$.
The inner SUM for Central is 0, and that 0 enters the outer MIN ---
identical to PER-WHEN.

\subsubsection*{Rewriting: can PER-WHEN replicate WHEN-PER?}

Under WHEN-PER, Central is excluded from the outer MIN entirely. Under
PER-WHEN, Central always participates. To replicate WHEN-PER, we would
need to remove Central from the outer aggregate.

The \textbf{coefficient trick} cannot do this. It operates at row level,
which controls the \emph{inner} SUM. Central's inner SUM is already 0 (all
coefficients zeroed out) --- but that 0 still enters the outer MIN. There
is no row-level operation that removes a group from an across-group
aggregate.

The \textbf{existence constraint} fails for a different reason. Adding
\texttt{SUM(CASE WHEN p.is\_organic THEN 1 ELSE 0 END) >= 1 PER p.region}
makes the model infeasible ($0 \geq 1$ for Central). But infeasibility is
the wrong outcome --- Central is not violating a requirement, it is simply
irrelevant to the organic revenue objective. The desired behavior is to
\emph{exclude} the group, not reject the model.

PER-WHEN cannot replicate WHEN-PER when the outer aggregate is MIN.
For outer SUM, MAX, and AVG, the orderings already agree, so no rewriting
is needed.



\end{document}